%% Chapter 9, Section 9.1 — Bayesian Filtering
%% A First Course in Monte Carlo Methods
%% Sanz-Alonso & Al-Ghattas

\section{Bayesian Filtering}
\label{sec:9.1}

Hidden Markov models arise in applications where a latent or ``hidden'' Markov process
$\{X_j\}_{j=0}^{\infty}$ needs to be estimated based on an \emph{observed} process
$\{Y_j\}_{j=1}^{J}$. For concreteness, we assume that the hidden and observed processes
are specified as follows:
\begin{alignat*}{3}
  \text{(Initial condition)}  &\quad& X_0 &\sim f_0, \\
  \text{(Dynamics model)}     &\quad& X_{j+1} &= a(X_j) + \xi_{j+1}, &\quad
    \xi_{j+1} &\sim \Normal(0, \Sigma),\quad j = 0, 1, \ldots, \\
  \text{(Observation model)}  &\quad& Y_{j+1} &= b(X_{j+1}) + \eta_{j+1}, &\quad
    \eta_{j+1} &\sim \Normal(0, \Gamma),\quad j = 0, 1, \ldots,
\end{alignat*}
where $X_0$, $\{\xi_j\}$, and $\{\eta_j\}$ are all independent; see Figure~\ref{fig:9.1}
for a representation of the resulting dependence structure of the hidden and observed
processes. The initial distribution $f_0$, the dynamics map $a : \R^d \to \R^d$, the
observation map $b : \R^d \to \R^k$, and the positive definite covariances $\Sigma$ and
$\Gamma$ are all assumed to be known. We assume Gaussianity of $\xi_j$ and $\eta_j$ for
pedagogical reasons, as it leads to concrete and familiar expressions for the Markov
kernel and likelihood function implicitly defined, respectively, by the dynamics and
observation models.

\begin{figure}[htbp]
  \centering
  \includegraphics[width=0.65\textwidth]{figures/fig_09_1}
  \caption{\textit{Graphical representation of the assumed dependence structure.}}
  \label{fig:9.1}
\end{figure}

We are interested in estimating $X_j$ given the observations $Y_{1:j} :=
\{Y_i\}_{i=1}^{j}$ available up to time $j$. In the Bayesian framework, inference is
carried out via the posterior or \emph{filtering distribution} of $X_j | Y_{1:j}$,
denoted by $f_j$. We will study two sequential Monte Carlo algorithms to approximate
$f_j$: the bootstrap particle filter and the optimal particle filter. An important feature
of these algorithms is that they are \emph{sequential}: at time $j+1$ they approximate
the distribution $f_{j+1}$ using an approximation of $f_j$ and the new observation
$Y_{j+1}$. To obtain a particle approximation of the filtering distribution, each
algorithm relies on a different conceptual decomposition of the filtering step:

\begin{enumerate}
  \item Bootstrap decomposition: in the bootstrap particle filter, particles are propagated
    through the dynamics model, and then Bayes' formula is applied to assimilate the new
    observation, leading to the following decomposition of the filtering step:
    $X_j | Y_{1:j} \longrightarrow X_{j+1} | Y_{1:j} \longrightarrow X_{j+1} | Y_{1:j+1}$.

  \item Optimal decomposition: in the optimal particle filter, Bayes' formula is applied
    to particles before they are propagated through the dynamics, leading to the following
    decomposition of the filtering step:
    $X_j | Y_{1:j} \longrightarrow X_j | Y_{1:j+1} \longrightarrow X_{j+1} | Y_{1:j+1}$.
\end{enumerate}

Throughout the chapter, we will use the notation
\[
  \pi(x) = \frac{1}{N} \sum_{n=1}^{N} \delta\bigl(x - X^{(n)}\bigr)
\]
to represent the empirical measure defined by a sample $\{X^{(n)}\}_{n=1}^{N}$, so that,
for sufficiently smooth test function $h$,
\[
  \calI_\pi[h] = \mathbb{E}_\pi[h] = \frac{1}{N} \sum_{n=1}^{N} h\bigl(X^{(n)}\bigr).
\]

We informally treat a distribution of this form as a p.d.f. Likewise, given normalized
weights $\{w^{(n)}\}_{n=1}^{N}$ and a sample $\{X^{(n)}\}_{n=1}^{N}$, we will use the
notation
\[
  \pi_w(x) = \sum_{n=1}^{N} w^{(n)} \delta\bigl(x - X^{(n)}\bigr)
\]
to represent the weighted empirical measure defined by the requirement that, for
sufficiently smooth test function $h$,
\[
  \calI_{\pi_w}[h] = \mathbb{E}_{\pi_w}[h] = \sum_{n=1}^{N} w^{(n)} h\bigl(X^{(n)}\bigr).
\]

Finally, for positive definite matrix $A$ and vector $v$, we will denote by
$|v|_A^2 := v^T A^{-1} v$ the Mahalanobis norm.
